\documentclass[11pt,a4paper]{article}
\usepackage[T1]{fontenc}
\usepackage[utf8]{inputenc}
\usepackage{amsmath,amssymb,amsthm}
\usepackage[margin=1in]{geometry}
\usepackage[colorlinks=true,linkcolor=blue,urlcolor=blue]{hyperref}
\theoremstyle{plain}
\newtheorem{theorem}{Theorem}
\newtheorem{lemma}[theorem]{Lemma}
\newtheorem{proposition}[theorem]{Proposition}
\newtheorem{corollary}[theorem]{Corollary}
\theoremstyle{definition}
\newtheorem{remark}[theorem]{Remark}

\title{The empirical recurrence for OEIS A154140, proved from an exact model}
\author{Adrian Perez Fontelles\\ \small Independent researcher}
\date{14 September 2026}
\begin{document}
\maketitle

\begin{abstract}
OEIS A154140 carries an empirical recurrence of order $5$. It is true, and it is
decidable rather than empirical. The entry's sequence is given exactly by the solutions of a Pell equation, in their finitely many orbits, from
which $a$ provably satisfies a MONIC linear recurrence of order $S=6$, derived below and not
assumed. The residual of the conjectured recurrence satisfies the same one, so whether it
annihilates $a$ is settled by a finite exact computation. No transfer matrix is involved: the sequence is the ordered list of solutions of a Pell equation, and what makes it $C$-finite is the automorph of the quadratic form.
\end{abstract}

\noindent\small 2020 Mathematics Subject Classification. 11D09, 11B37, 11D45.\normalsize

\section{The conjecture}

OEIS A154140 is ``Indices k such that 6 plus the k-th triangular number is a perfect square.''. It has offset $1$ and begins
\[
2, 4, 19, 29, 114, 172, 667, 1005,\ \dots
\]
The entry states, as a conjecture contributed by its author and never marked settled:
\begin{quote}
Conjectures: (Start) a(n) = +a(n-1) +6*a(n-2) -6*a(n-3) -a(n-4) +a(n-5).
\end{quote}
As of the ``Last modified'' line on the live entry (Dec 23 2024 14:53:42, revision 29)
this is still recorded as empirical, and nothing on the entry records it as proved.

\section{The count is a Pell equation}

$k(k+1)/2 + 6$ is a perfect square exactly when $8$ times it is, and completing the square
turns that into
\[
(2k+1)^2-2(2m)^2\;=\;1-8\cdot6\;=\;-47,
\]
so writing $X=2k+1$ and $Y=2m$ the admissible $k$ are exactly the solutions of $X^2-2Y^2=-47$
with $X$ odd and positive and $Y$ even and nonnegative. The entry's sequence is that solution
set, listed by increasing $k$.

The form $x^2-2y^2$ has the automorph coming from the unit $3+2\sqrt2$,
\[
T(X,Y)=(3X+4Y,\;2X+3Y),
\]
which carries solutions to solutions; and every solution is $T^{j}$ of a FUNDAMENTAL one,
meaning one whose preimage $T^{-1}(X,Y)=(3X-4Y,-2X+3Y)$ leaves the region. A finite search
finds them -- there are $2$ here -- and the search is checked rather than bounded by a quoted
theorem: it is run to a limit and then re-run to four times that limit, and the engine refuses
unless the two agree.

\section{The bound}

On the region $X\ge1$, $Y\ge0$ one has $Y=\sqrt{(X^2--47)/2}$, an increasing function of $X$,
so $T(X)=3X+4Y$ is strictly increasing: $T$ PRESERVES THE ORDER of solutions. Hence if the
sorted list of solutions has its $(n+2)$-th entry equal to $T$ of its $n$-th for $2$
consecutive $n$, it does so for every later $n$ -- the $2$ orbits interleave in a fixed cyclic
order from there on. That happens from index $n_0=0$, checked directly.

Within one orbit $X_{j+2}=6X_{j+1}-X_j$, since $3+2\sqrt2$ and $3-2\sqrt2$ are the roots of
$t^2-6t+1$; in terms of $k=(X-1)/2$ that reads $k_{j+2}=6k_{j+1}-k_j+2$. Combining with the
interleaving,
\[
a(n+2\cdot2)\;=\;6\,a(n+2)-a(n)+2 \qquad (n\ge n_0),
\]
so $(z-1)(z^{2\cdot2}-6z^{2}+1)$ annihilates $a$, monic of order $S=6$ once the
transient is allowed for. Nothing in that bound is estimated: $2$ is the number of fundamental
solutions and $n_0$ is where the interleaving is observed to settle, which the order-preservation
of $T$ then makes permanent.

\section{The proof}

\begin{lemma}\label{lem:crit}
Let $A(z)=z^S-\sum_{j<S}\alpha_jz^{j}$ be the monic annihilator derived above and
$q(z)=z^{5}-\sum_ic_iz^{5-i}$ the characteristic polynomial of the conjectured
recurrence. The residual $r(n)=a(n)-\sum_ic_ia(n-i)$ is a linear combination of shifts of $a$,
so $A$ annihilates $r$ as well. Since $A(0)\ne0$ the recurrence it gives runs in both
directions, and $S$ consecutive zeros of $r$ force $r$ to vanish at every later index.
\end{lemma}

\begin{proof}
$A$ annihilates $a$, hence every shift of $a$, hence any linear combination of shifts; that is
$r$. A monic recurrence with nonzero constant term determines each term from the $S$ before it
and each from the $S$ after it, so a run of $S$ zeros propagates.
\end{proof}

\begin{theorem}
\[
a(n)\;=\;a(n-1) + 6\,a(n-2) - 6\,a(n-3) - a(n-4) + a(n-5)
\]
for every $n>1$.
\end{theorem}

\begin{proof}
The terms $a(0),\dots$ were computed exactly in integer arithmetic from the model, extended by
$A$ where the model itself was not evaluated, and $r(n)$ formed for each index. Those integers
vanish from $n=1+1$ onwards, and the run exceeds $S+5$, which by
Lemma~\ref{lem:crit} settles every larger index at once.
\end{proof}

\section{Verification}

Three checks, each independent of the algebra above.

First, the model reproduces all $30$ terms the entry publishes, exactly, in integer
arithmetic. This is what ties the model to the entry: it is built by reading the entry's
English, and a misreading gives different counts.

Second, the conjectured recurrence was evaluated directly on those published terms, with no
model involved, and holds wherever the proved range and the data overlap.

Third, the derived annihilator $A$ was itself tested on terms it had not been given: the model
was evaluated beyond the $S$ values that determine it and $A$ reproduced them. A bound that is
too small fails this test, which is why it is run.

\begin{thebibliography}{9}
\bibitem{oeis} The OEIS Foundation, \emph{The On-Line Encyclopedia of Integer
Sequences}, \url{https://oeis.org}, 2026. Sequence A154140.
\bibitem{beckrobins} M.~Beck and S.~Robins, \emph{Computing the Continuous Discretely},
2nd ed., Springer, 2015. (Ehrhart quasi-polynomials and their periods.)
\bibitem{stanley} R.~P.~Stanley, \emph{Enumerative Combinatorics}, Volume 1, 2nd ed.,
Cambridge University Press, 2011.
\bibitem{ds} F.~Diamond and J.~Shurman, \emph{A First Course in Modular Forms},
Springer GTM 228, 2005.
\end{thebibliography}

\end{document}
